\section*{Chapitre \ref{ch:in:machine-learning-and-data-engineer} --- Apprentissage automatique et ingénierie des données}

\begin{solution}{exo:in:machine-learning-and-data-engineer:1}
$1\,051/136=7.73$~; sur quatre ans, $7.73^{1/4}-1=66.7~\%$ par an.
\end{solution}

\begin{solution}{exo:in:machine-learning-and-data-engineer:2}
$136/4\,422=3.1~\%$ pour l'exercice 2021 et $1\,051/8\,808=11.9~\%$ pour l'exercice 2025.
\end{solution}

\begin{solution}{exo:in:machine-learning-and-data-engineer:3}
La finance et l'assurance emploient davantage (46\,730 contre 32\,410)~; le secteur de l'information verse la médiane
la plus élevée (141\,440~\$ contre 124\,770~\$).
\end{solution}

\begin{solution}{exo:in:machine-learning-and-data-engineer:4}
$\sqrt{1/2\,618+1/5\,727}/4=0.0059$, contre 0.023~: l'erreur type d'un effectif logarithmique vaut environ $1/\sqrt n$, et
les effectifs des ingénieurs logiciels sont vingt à quarante fois plus grands.
\end{solution}

\begin{solution}{exo:in:machine-learning-and-data-engineer:5}
Les dépôts de l'exercice 2021 ont été codés selon la classification de 2010, qui ne comportait pas de profession \omterm{def:in:machine-learning-and-data-engineer:roles}{data scientist} (ces postes étaient déclarés
comme statisticiens, analystes en recherche opérationnelle, informaticiens ou sous d'autres codes). La famille de titres est définie par les
intitulés de poste, qui existent dans les deux exercices.
\end{solution}

\begin{solution}{exo:in:machine-learning-and-data-engineer:6}
Les dépôts indiquent que les offres de \omterm{def:in:machine-learning-and-data-engineer:roles}{data scientist} dans les banques sont inférieures à celles des développeurs logiciels au même niveau, et que
les chercheurs des fonds systématiques reçoivent des offres bien plus élevées (chapitre~17)~; la croissance des dépôts montre que le poste bancaire est recherché. Ils
ne disent pas ce que rapporte l'un ou l'autre poste en primes, quel travail ferait le \omterm{def:in:machine-learning-and-data-engineer:roles}{data scientist} de la banque, ni en quoi les deux carrières divergent.
\end{solution}

\begin{solution}{exo:in:machine-learning-and-data-engineer:7}
La variance de chaque effectif vaut dix fois sa moyenne, si bien que l'erreur type est multipliée par $\sqrt{10}$~: 0.072 au lieu de 0.023. La
croissance annuelle reste à 66.7~\%, mais son intervalle s'élargit de 44.8~\% à 92.0~\%.
\end{solution}

\begin{solution}{exo:in:machine-learning-and-data-engineer:8}
Neuf dépôts sur dix d'apprentissage automatique en 2025 proviennent de banques, et non de sociétés de trading~; les dépôts
recensent des demandes de visa pour de nouvelles embauches, et non l'emploi~; la stagnation des dépôts de chercheurs n'est pas une baisse~; et les chercheurs
d'une société de trading utilisent l'apprentissage automatique sous leur propre intitulé, de sorte qu'un changement de méthode n'entraîne pas forcément un changement d'intitulé.
\end{solution}

\begin{solution}{pb:in:machine-learning-and-data-engineer:1}
\begin{enumerate}
\item Comme dans la définition du chapitre.
\item Le problème~: un rapport signal sur bruit faible, peu d'observations indépendantes, une non-stationnarité et un coût monétaire de l'erreur.
\item Les systèmes qui entraînent, servent et surveillent les modèles~; le contrat de ce que doit contenir un paquet de
modèle entre ceux qui le construisent et ceux qui l'exploitent (livre~12).
\item Le chemin depuis les fichiers des fournisseurs et des plateformes jusqu'aux requêtes des chercheurs~: capture, stockage, versions à l'instant des données,
feature stores, contrôles.
\item Il n'existe pas de profession d'\omterm{def:in:machine-learning-and-data-engineer:roles}{ingénieur de données}~; O*NET classe l'intitulé parmi les architectes de bases de
données, et la plupart des postes sont déclarés comme développeurs logiciels.
\item 136 et 1\,051.
\item Des effectifs de Poisson à moyenne log-linéaire~; $g=\log(n_1/n_0)/4$, erreur type $\sqrt{1/n_0+1/n_1}/4$.
\item 66.7~\% (59.4~\% à 74.3~\%)~; 21.6~\% (20.2~\% à 23.0~\%)~; 0.2~\% ($-1.7~\%$ à 2.1~\%).
\item Elle suppose des dépôts indépendants~; des lots venant de quelques employeurs élargissent l'incertitude réelle.
\item Banques~: 920 sur 1\,051.
\item $-\$17\,000$ au total ($-\$20\,000$ à $-\$15\,000$)~; $-\$20\,500$, $-\$18\,300$, $-\$19\,000$ aux niveaux I à III~;
$+\$2\,220$ au niveau IV.
\item Des médianes proches de celles des chercheurs et ingénieurs de ces mêmes entreprises~: de 175\,000~\$ à 193\,000~\$, à partir de dix à douze dépôts chacune.
\item Services professionnels et techniques, finance et assurance, et information~; dans la finance, surtout la banque et
l'assurance.
\item Les dépôts ne le disent pas~; le poste peut servir des activités moins bien payées que l'ingénierie, ou être déclaré côté analyste
de l'intitulé.
\item Primes et actions.
\item 66.7~\% par an (59.4~\% à 74.3~\%)~; environ 19\,000~\$ de moins que les développeurs logiciels aux niveaux I à III, à égalité au niveau IV.
\item Surtout pas~: la plus grande part se trouve dans les banques.
\item Le master de statistique prépare à la data science et à l'apprentissage automatique intégré à la recherche~; le
master de mathématiques financières prépare aux postes de quant et de chercheur en banque (chapitre~29).
\item Quelle activité le poste sert, ce qu'il livre, qui possède les modèles, et comment sa rémunération se compare à celle des ingénieurs.
\item Les recrutements sur les intitulés d'apprentissage automatique et de données dans la finance ont crû d'environ deux tiers par an jusqu'à l'exercice 2025, surtout dans les banques
et à des niveaux de salaire inférieurs à ceux des ingénieurs. Il faut choisir le master selon le poste visé~: la recherche de trading recrute encore sous un intitulé de recherche,
et utilise l'apprentissage automatique comme méthode.
\end{enumerate}
\end{solution}

\begin{solution}{iq:in:machine-learning-and-data-engineer:1}
Mélanger les données revient à placer des observations futures dans les blocs d'entraînement des points de test passés, et laisse fuir des étiquettes qui se chevauchent
d'un bloc à l'autre~; il faut utiliser une découpe glissante dans le temps, ou une découpe purgée et avec embargo (livre~12, chapitre~3).
\iqlookfor{fuite à travers le temps et étiquettes qui se chevauchent.}
\end{solution}

\begin{solution}{iq:in:machine-learning-and-data-engineer:2}
Une variable calculée différemment en production~; un regard vers l'avenir dans les données d'entraînement~; un changement de fournisseur de données ou de marché~;
un surapprentissage sur l'ensemble de validation~; et une dérive de la population. Il faut vérifier d'abord ce qui est peu coûteux et certain~: les variables, puis les données, puis
le modèle.
\iqlookfor{écart entraînement--service avant la dérive.}
\end{solution}

\begin{solution}{iq:in:machine-learning-and-data-engineer:3}
Une seule définition dans un feature store utilisé par les deux, des jointures à l'instant des données pour l'entraînement, et une comparaison en direct des valeurs servies
avec celles recalculées.
\iqlookfor{une seule définition et un contrôle.}
\end{solution}

\begin{solution}{iq:in:machine-learning-and-data-engineer:4}
Des contrôles de schéma et de distribution à l'arrivée (types, plages, taux de valeurs nulles, quantiles comparés à l'historique), un rapprochement avec
une seconde source lorsqu'elle existe, et des alertes qui arrêtent les traitements en aval.
\iqlookfor{contrats de données et surveillance des distributions.}
\end{solution}

\begin{solution}{iq:in:machine-learning-and-data-engineer:5}
Seulement si le gain résiste aux coûts et au temps~: il faut le tester hors échantillon sur plusieurs périodes, mesurer sa valeur
monétaire après coûts, et peser la maintenance et le risque du modèle plus complexe~; un gain de 0.3~\% en précision peut n'être que du bruit.
\iqlookfor{valeur économique et stabilité, non précision.}
\end{solution}

\begin{solution}{iq:in:machine-learning-and-data-engineer:6}
Partitionner par date et par symbole dans un format en colonnes avec compression, tenir un index des plages temporelles par fichier, stocker les versions brutes et
normalisées, et servir les requêtes par un moteur qui élague les partitions (livre~15, chapitres~3 et~4).
\iqlookfor{partitionnement selon le mode d'accès et stockage en colonnes.}
\end{solution}
